\documentclass[11pt]{article}
\usepackage[margin=1in]{geometry}
\usepackage{amsmath,amssymb,amsthm}
\usepackage{graphicx}
\usepackage{hyperref}

\newtheorem{theorem}{Theorem}
\newtheorem{lemma}{Lemma}
\newtheorem{corollary}{Corollary}

\newcommand{\R}{\mathbb R}
\newcommand{\N}{\mathbb N}
\newcommand{\Id}{\mathrm{Id}}
\newcommand{\spanop}{\mathrm{span}}
\newcommand{\wstar}{w^*}

\title{The Codimension-Two Full Interval for Non-Attaining Projections in \(C[0,1]\)}
\author{Automatic Math Discovery Smoke Test}
\date{June 29, 2026}

\begin{document}
\maketitle

\section*{Source and Extracted Question}

The source is Tomasz Kania and Grzegorz Lewicki, \emph{Finite-codimensional
subspaces of Daugavet spaces: projection constants and minimal projections},
arXiv:2604.10771.  Question 6.2 asks whether, for each fixed \(n\in\N\), every
value in the interval
\[
  (2,1+\lambda_n]
\]
occurs as the projection constant of a codimension-\(n\) subspace \(Y\) of the
real space \(C[0,1]\) for which no minimal projection exists.  Here
\(\lambda_n\) is the maximal absolute projection constant among real
\(n\)-dimensional Banach spaces.

\begin{center}
\includegraphics[width=\textwidth]{figures/open_problem_crop.png}
\end{center}

The present packet settles this question for \(n=2\).  We use the classical
value \(\lambda_2=4/3\), attained by the regular hexagon norm.

\section*{Proof Intuition}

Kania--Lewicki reduce finite-codimensional subspaces of a Daugavet space to
finite-dimensional subspaces of the dual.  In the codimension-two case it is
therefore enough to build two-dimensional subspaces \(W\subset M[0,1]\) with
the desired projection constants and with no \(\wstar\)-continuous minimal
projection.

The useful models are the hyperplanes
\[
  V_t=\{x\in \ell_1^3:x_1+x_2+t x_3=0\},\qquad 0<t\leq 1.
\]
Their projection constants fill the whole interval \((1,4/3]\).  The new point
is that no duplication-stability argument is needed.  The unique minimal
projection onto \(V_t\) forces the interval averages of any
\(\wstar\)-continuous minimal projection to land on prescribed exposed faces of
the coefficient ball
\[
  |a+tb|+|a|+|b|\leq r(t).
\]
On two adjacent intervals these faces are disjoint.  Continuous coefficient
functions cannot approach one face from the left and a disjoint face from the
right at the common endpoint.

\section*{Finite-Dimensional Model}

For \(0<t\leq 1\), put
\[
  D=1+t+t^2,\qquad
  r(t)=\frac{(1+t)^2}{D},
\]
and define
\[
  V_t=\ker s_t\subset \ell_1^3,\qquad
  s_t(x_1,x_2,x_3)=x_1+x_2+t x_3.
\]

\begin{lemma}
\label{lem:finite}
The space \(V_t\) has a unique minimal projection in \(\ell_1^3\).  It is
\[
  P_t=\Id-u_t\otimes s_t,\qquad
  u_t=\left(\alpha,\alpha,\frac{t}{D}\right),
  \quad
  \alpha=\frac{1+t}{2D},
\]
and
\[
  \lambda(V_t,\ell_1^3)=\|P_t\|=r(t).
\]
\end{lemma}

\begin{proof}
Every projection from \(\ell_1^3\) onto \(V_t\) has the form
\[
  P_u=\Id-u\otimes s_t,\qquad s_t(u)=1.
\]
Let \(e_1,e_2,e_3\) be the coordinate basis and set
\[
  A_1=\|e_1-u\|_1,\quad
  A_2=\|e_2-u\|_1,\quad
  A_3=\|e_3-tu\|_1.
\]
Then \(\|P_u\|=\max_j A_j\).

Consider the three sign functionals
\[
  \phi_1=(1,-1,-1),\qquad
  \phi_2=(-1,1,-1),\qquad
  \phi_3=(-1,-1,1)
\]
and the positive weights
\[
  \omega_1=\omega_2=\frac{t(1+t)}{2D},\qquad
  \omega_3=\frac{1}{D}.
\]
They sum to \(1\), and a direct calculation gives
\[
  \omega_1\phi_1+\omega_2\phi_2+t\omega_3\phi_3
  =-\frac{t}{D}(1,1,t).
\]
Since \(\|\phi_j\|_{\infty}=1\), for every admissible \(u\),
\[
\begin{aligned}
  \|P_u\|
  &\geq \sum_{j=1}^3 \omega_j A_j                                      \\
  &\geq
  \omega_1\phi_1(e_1-u)+\omega_2\phi_2(e_2-u)+\omega_3\phi_3(e_3-tu)       \\
  &=1+\frac{t}{D}s_t(u)
  =1+\frac{t}{D}
  =\frac{(1+t)^2}{D}
  =r(t).
\end{aligned}
\]

For the stated \(u_t\), one has \(s_t(u_t)=1\), and the three columns of
\(P_t\) have signs
\[
  (+,-,-),\qquad (-,+,-),\qquad (-,-,+),
\]
respectively.  Hence the above inequalities are equalities, so
\(\|P_t\|=r(t)\).

It remains to prove uniqueness.  If \(P_u\) is minimal, equality holds in the
weighted estimate.  Since all weights are positive and each \(A_j\leq r(t)\),
we must have \(A_j=r(t)\) for each \(j\), and equality must hold in the
corresponding sign estimates.  In particular,
\[
  \phi_1(e_1-u)=r(t),\qquad \phi_2(e_2-u)=r(t).
\]
These two equations are
\[
  1-u_1+u_2+u_3=r(t),\qquad
  1+u_1-u_2+u_3=r(t),
\]
so \(u_1=u_2\).  The first equation then gives
\[
  u_3=r(t)-1=\frac{t}{D}.
\]
Together with \(s_t(u)=1\), this yields
\[
  u_1=u_2=\frac{1-tu_3}{2}
  =\frac{1+t}{2D}.
\]
Thus \(u=u_t\), and the minimal projection is unique.
\end{proof}

\section*{A Copy in \(M[0,1]\)}

Let \(I_1=(0,1/3)\), \(I_2=(1/3,2/3)\), and \(I_3=(2/3,1)\), and let
\[
  \mu_j=3\,\mathbf 1_{I_j}\,dt\qquad (j=1,2,3).
\]
The span \(E=\spanop\{\mu_1,\mu_2,\mu_3\}\subset M[0,1]\) is isometric to
\(\ell_1^3\).  Define
\[
  W_t=
  \left\{
    z_1\mu_1+z_2\mu_2+z_3\mu_3:
    z_1+z_2+t z_3=0
  \right\}.
\]

\begin{lemma}
\label{lem:measure}
For every \(0<t\leq 1\),
\[
  \lambda(W_t,M[0,1])=r(t),
\]
but \(W_t\) admits no \(\wstar\)-continuous minimal projection from \(M[0,1]\)
onto \(W_t\).
\end{lemma}

\begin{proof}
The equality of projection constants is immediate from the finite-dimensional
model.  Restricting any projection \(M[0,1]\to W_t\) to \(E\) gives a
projection \(E\to W_t\), so
\[
  \lambda(W_t,M[0,1])\geq \lambda(V_t,\ell_1^3)=r(t).
\]
Conversely, the map
\[
  \nu\longmapsto \sum_{j=1}^3 \nu(I_j)\mu_j
\]
is a contractive projection from \(M[0,1]\) onto \(E\).  Composing it with the
minimal projection \(E\to W_t\) gives a projection of norm \(r(t)\).  Hence the
projection constant is exactly \(r(t)\).

Suppose, toward a contradiction, that \(Q:M[0,1]\to W_t\) is a
\(\wstar\)-continuous projection with \(\|Q\|=r(t)\).  Use the basis
\[
  v_1=\mu_1-\mu_2,\qquad v_2=t\mu_1-\mu_3
\]
of \(W_t\).  Since \(Q\) has finite rank and is \(\wstar\)-continuous, there
are continuous real functions \(g,h\in C[0,1]\) such that
\[
  Q\nu=\left(\int g\,d\nu\right)v_1+\left(\int h\,d\nu\right)v_2
  \qquad (\nu\in M[0,1]).
\]
For \(a,b\in\R\), put
\[
  N_t(a,b)=|a+tb|+|a|+|b|.
\]
Then
\[
  \|a v_1+b v_2\|=N_t(a,b),
\]
and applying \(Q\) to point masses gives
\[
  N_t(g(x),h(x))=\|Q\delta_x\|\leq r(t)\qquad (x\in[0,1]).
\]

Let
\[
  c_j=\left(\int g\,d\mu_j,\int h\,d\mu_j\right)
  \qquad (j=1,2,3).
\]
The restriction of \(Q\) to \(E\) is a projection \(E\to W_t\) of norm at most
\(r(t)\).  By Lemma~\ref{lem:finite}, it is the unique minimal projection.
Therefore the first two coefficient columns are
\[
  c_1=\left(\alpha,\frac{t}{D}\right),
  \qquad
  c_2=\left(\alpha-1,\frac{t}{D}\right),
  \qquad
  \alpha=\frac{1+t}{2D}.
\]

Now define the two linear functionals
\[
  L_1(a,b)=2a+(1+t)b,\qquad
  L_2(a,b)=-2a+(1-t)b.
\]
For all \(a,b\in\R\), \(L_1(a,b)\leq N_t(a,b)\) and
\(L_2(a,b)\leq N_t(a,b)\).  Moreover,
\[
  L_1(c_1)=r(t),\qquad L_2(c_2)=r(t).
\]
Since \(c_1\) and \(c_2\) are the \(\mu_j\)-averages of the continuous
coefficient curve \(x\mapsto (g(x),h(x))\), while \(L_i(g(x),h(x))\leq r(t)\)
pointwise, we get
\[
  L_1(g(x),h(x))=r(t)\quad\text{for a.e. }x\in I_1,
\]
and
\[
  L_2(g(x),h(x))=r(t)\quad\text{for a.e. }x\in I_2.
\]
By continuity and closedness, at the common endpoint \(x=1/3\) the point
\[
  (a,b)=(g(1/3),h(1/3))
\]
must satisfy
\[
  N_t(a,b)\leq r(t),\qquad L_1(a,b)=r(t),\qquad L_2(a,b)=r(t).
\]
The last two equalities imply \(b=r(t)\).  Hence
\[
  N_t(a,b)=|a+t r(t)|+|a|+r(t)\leq r(t),
\]
which forces \(a=0\) and \(a=-t r(t)\).  This is impossible because
\(t>0\) and \(r(t)>0\).  Thus no \(\wstar\)-continuous minimal projection
exists.
\end{proof}

\section*{Transfer to \(C[0,1]\)}

\begin{theorem}
\label{thm:main}
For every
\[
  2<\Lambda\leq \frac{7}{3},
\]
there is a codimension-\(2\) subspace \(Y\) of the real space \(C[0,1]\) such
that
\[
  \lambda(Y,C[0,1])=\Lambda,
\]
and no minimal projection from \(C[0,1]\) onto \(Y\) exists.
\end{theorem}

\begin{proof}
The function
\[
  r(t)=\frac{(1+t)^2}{1+t+t^2}
\]
is continuous and strictly increasing on \((0,1]\), with
\[
  \lim_{t\downarrow 0}r(t)=1,\qquad r(1)=\frac43.
\]
Thus for every \(\Lambda\in(2,7/3]\) there exists \(t\in(0,1]\) such that
\[
  \Lambda=1+r(t).
\]

Let \(W_t\subset M[0,1]=C[0,1]^*\) be the two-dimensional space from
Lemma~\ref{lem:measure}, and set
\[
  Y=W_t^\perp\subset C[0,1].
\]
The source paper's Daugavet duality theorem says that for a finite-codimensional
subspace \(Y=W^\perp\) of a Daugavet space \(X\),
\[
  \lambda(Y,X)=1+\lambda(W,X^*),
\]
and that a minimal projection onto \(Y\) exists if and only if there is a
\(\wstar\)-continuous minimal projection from \(X^*\) onto \(W\).  Applying
this with \(X=C[0,1]\) and \(W=W_t\), Lemma~\ref{lem:measure} gives
\[
  \lambda(Y,C[0,1])=1+r(t)=\Lambda,
\]
and no minimal projection onto \(Y\) exists.  Since \(\dim W_t=2\), the
subspace \(Y\) has codimension \(2\).
\end{proof}

\begin{corollary}
Question 6.2 of arXiv:2604.10771 has an affirmative answer for \(n=2\).
\end{corollary}

\begin{proof}
The classical two-dimensional maximal absolute projection constant is
\(\lambda_2=4/3\).  Therefore the interval in Question 6.2 is exactly
\((2,7/3]\), and Theorem~\ref{thm:main} realizes every value in it with no
minimal projection.
\end{proof}

\section*{Novelty Check and Limitations}

The cheap run indexes contained only the earlier failed attempt under the slug
\emph{codim2 endpoint regular-hexagon route}.  That attempt identified the
family \(V_t\), but stopped because block duplicates do not have unique minimal
projections.  The present proof bypasses duplication entirely by using the
equality case of the original three-coordinate projection problem and an
exposed-face continuity obstruction.

Bounded web searches on 2026-06-29 for the arXiv id, the paper title,
codimension-two/full-interval phrases, and the \(\ker(1,1,t)\) projection
constant route found no later explicit solution.  This novelty check is not a
formal literature review.

The theorem settles the fixed-codimension question only for \(n=2\).  It does
not address the remaining cases \(n\geq 3\).

\begin{thebibliography}{9}

\bibitem{KL}
T. Kania and G. Lewicki,
\emph{Finite-codimensional subspaces of Daugavet spaces: projection constants
and minimal projections},
arXiv:2604.10771, 2026.

\bibitem{KTJ}
H. K\"onig and N. Tomczak-Jaegermann,
\emph{Bounds for projection constants and \(1\)-summing norms},
Trans. Amer. Math. Soc. \textbf{320} (1990), 799--823.

\end{thebibliography}

\end{document}
